\documentclass[11pt]{article}
\usepackage[a4paper,margin=1in]{geometry}
\usepackage{fontspec}
\usepackage[boldfont=false]{xeCJK}
\setCJKmainfont{FandolSong-Regular.otf}
\usepackage{amsmath}
\usepackage{amssymb}
\usepackage{array}
\usepackage{booktabs}
\usepackage{graphicx}
\usepackage{adjustbox}
\usepackage{enumitem}
\setlist[itemize]{topsep=2pt,partopsep=0pt,parsep=0pt,itemsep=2pt,leftmargin=1.4em}
\setlist[enumerate]{topsep=2pt,partopsep=0pt,parsep=0pt,itemsep=2pt,leftmargin=1.6em}
\usepackage{parskip}
\usepackage{fvextra}
\fvset{breaklines=true,breakanywhere=true,fontsize=\small}
\setlength{\parindent}{0pt}

\begin{document}

\begin{center}
  {\LARGE\bfseries Photosynthesis}\\[0.35em]
  {\large A Level Biology}
\end{center}
\vspace{0.6em}

\section*{The chloroplast}

\textbf{Photosynthesis} 光合作用 happens inside the \textbf{chloroplast} 叶绿体. Its structure suits its two stages:

\begin{itemize}
\item inside are stacks of flat sacs called \textbf{thylakoids} 类囊体. A stack of thylakoids is a \textbf{granum} (plural \textbf{grana} 基粒). The thylakoid membranes hold the light-trapping pigments and are the site of the first stage.

\item the fluid around the thylakoids is the \textbf{stroma} 基质, the site of the second stage.

\end{itemize}

\section*{The two stages of photosynthesis}

Photosynthesis has two linked stages:

\begin{enumerate}
\item the \textbf{light-dependent stage} 光反应阶段 happens in the thylakoids. It uses light energy to make ATP and \textbf{reduced} 还原 NADP.

\item the \textbf{light-independent stage} 暗反应阶段, also called the \textbf{Calvin cycle} 卡尔文循环, happens in the stroma. It uses the ATP and reduced NADP from the first stage to build complex organic molecules from \textbf{carbon dioxide} 二氧化碳.

\end{enumerate}

\section*{Chloroplast pigments}

A \textbf{pigment} 色素 is a coloured substance that absorbs light. The thylakoids hold several pigments that work together to trap as much light as possible:

\begin{itemize}
\item \textbf{chlorophyll} 叶绿素 a and chlorophyll b (which absorb mainly red and blue light),

\item \textbf{carotene} 胡萝卜素 and \textbf{xanthophyll} 叶黄素 (which absorb other colours and pass the energy on).

\end{itemize}

We study them with two graphs. An \textbf{absorption spectrum} 吸收光谱 shows how much light each pigment absorbs at each wavelength. An \textbf{action spectrum} 作用光谱 shows how fast photosynthesis goes at each wavelength. The two graphs match closely, which shows the pigments drive photosynthesis.

You can separate the pigments by \textbf{chromatography} 色谱法: the pigments travel different distances up the paper. Each pigment is identified by its \textbf{Rf value} 比移值 (the distance the pigment moved divided by the distance the solvent moved).

\section*{The light-dependent stage}

In the thylakoids, light is used to make ATP. This is \textbf{photophosphorylation} 光合磷酸化, and it comes in two forms.

In \textbf{cyclic photophosphorylation} 循环光合磷酸化:

\begin{itemize}
\item only \textbf{photosystem} 光系统 I (PSI) is used.

\item \textbf{photoactivation} 光激活 of chlorophyll occurs (light boosts its \textbf{electrons} 电子 to a higher energy).

\item only ATP is made.

\end{itemize}

In \textbf{non-cyclic photophosphorylation} 非循环光合磷酸化:

\begin{itemize}
\item both photosystem I (PSI) and photosystem II (PSII) are used.

\item photoactivation of chlorophyll occurs.

\item the \textbf{oxygen-evolving complex} 放氧复合体 carries out the \textbf{photolysis} 光解 (splitting by light) of water, which releases \textbf{oxygen} 氧气.

\item both ATP and reduced NADP are made.

\end{itemize}

In both forms, energy is released in the same way:

\begin{enumerate}
\item energetic electrons pass along an \textbf{electron transport chain} 电子传递链, releasing energy as they go.

\item this energy is used to pump \textbf{protons} 质子 across the thylakoid membrane.

\item the protons flow back into the stroma through \textbf{ATP synthase} 合酶, and this flow provides the energy to make ATP.

\end{enumerate}

\section*{The Calvin cycle}

The light-independent stage builds sugars in three main steps:

\begin{enumerate}
\item \textbf{fixation} 固定 — the \textbf{enzyme} 酶 rubisco joins carbon dioxide to a 5-carbon molecule, RuBP (ribulose bisphosphate). This makes two molecules of a 3-carbon compound, GP (glycerate 3-phosphate).

\item \textbf{reduction} — GP is reduced to TP (triose phosphate) using reduced NADP and ATP from the light-dependent stage.

\item \textbf{regeneration} — most of the TP is used to \textbf{regenerate} 再生 the RuBP (using more ATP), so the cycle can keep running.

\end{enumerate}

Some TP leaves the cycle to make useful molecules. GP can be used to make some \textbf{amino acids} 氨基酸, and TP can be used to make \textbf{carbohydrates} 碳水化合物, \textbf{lipids} 脂质 and amino acids.

\section*{Limiting factors}

A \textbf{limiting factor} 限制因素 is the one in shortest supply that holds back the rate of photosynthesis. The three main ones are \textbf{light intensity} 光照强度, \textbf{carbon dioxide} \textbf{concentration} 浓度, and \textbf{temperature} 温度.

\begin{itemize}
\item raising light intensity speeds up photosynthesis, until some other factor becomes limiting.

\item raising carbon dioxide concentration speeds it up, until some other factor becomes limiting.

\item raising temperature speeds it up, but only to an optimum; too high and the enzymes denature.

\end{itemize}

You can measure the rate of the light-dependent stage with a redox \textbf{indicator} 指示剂 such as DCPIP or methylene blue and a suspension of chloroplasts: the dye loses its colour as the chloroplasts work, and you can test different light intensities or light \textbf{wavelengths} 波长. With a whole \textbf{aquatic plant} 水生植物 you can count the bubbles of oxygen given off to compare rates under different conditions.


\end{document}
